\chapter{资产配置：唯一的免费午餐}
\chaptermeta{21}

\begin{ledetext}
金融学里几乎所有东西都要你拿什么去换。只有一件事，收益不减，风险变小。它的条件比多数人以为的苛刻。
\end{ledetext}

\section{那顿饭长什么样}

上一章的结论是：长期回报来自承担风险。听着像是没得选，想要收益就得挨揍。

有一个例外。

假设有两个资产，长期收益都是 8\%，但它们涨跌的节奏对不上。各买一半，每年调回一半一半。看看会发生什么。

\begin{tablewrap}
\begin{xltabular}{\linewidth}{>{\raggedright\arraybackslash}p{0.20\linewidth}rrr}
\toprule
\thead{年份} & \thead{资产 A} & \thead{资产 B} & \thead{各一半，年末调回} \\
\midrule
\textbf{第 1 年} & +32\% & −9\% & +11.5\% \\
\textbf{第 2 年} & −18\% & +24\% & +3.0\% \\
\textbf{第 3 年} & +25\% & −5\% & +10.0\% \\
\textbf{第 4 年} & −11\% & +22\% & +5.5\% \\
\textbf{第 5 年} & +12\% & +8\% & +10.0\% \\
\textbf{\textbf{算术平均}} & 8.0\% & 8.0\% & 8.0\% \\
\textbf{\textbf{实际年化}} & 6.17\% & 7.15\% & 7.95\% \\
\textbf{\textbf{波动（标准差）}} & 19.6\% & 13.5\% & 3.2\% \\
\bottomrule
\end{xltabular}
\end{tablewrap}

10 万块放进去五年：A 变成 13.49 万，B 变成 14.12 万，各一半变成 \textbf{14.66 万}。

组合比它的两个成分都赚得多。这不是错觉，也不是我算错了。

原因在算术平均和实际年化之间那道缝。跌 18\% 之后要涨 22\% 才回本，波动越大，这道缝越宽，实际拿到手的越少于纸面上的平均数。把波动摁下去，缝就窄了，同样的 8\% 平均收益，你能带走更多。

\begin{pullquote}{}
收益没有减少，风险减少了。这是金融学里唯一一件近似免费的事。
\end{pullquote}

\begin{warnbox}{先说清楚这个例子的水分}
上面那两组数字是我故意挑的，它们几乎完全反向，\term{相关性}{correlation}\index{070@相关性}接近 −1。现实中找不到这么听话的一对。

股票和债券的相关性在过去二十年里从 −0.4 到 +0.6 之间来回跑，不是一个常数。所以真实世界的免费午餐没这么大份，效果通常是把波动降低两三成，而不是从 19.6\% 降到 3.2\%。

方向是对的，幅度要打折。任何一本把这个例子讲得像永动机的书，都跳过了这一段。

\end{warnbox}

\section{关键变量是相关性，不是数量}

免费午餐的开关只有一个：两样东西是不是一起动。

如果两个资产永远同涨同跌，各买一半和全买一个，波动一模一样，什么都没省下。买十只、买五十只，同样什么都没省下。

\begin{mythbox}{常见误解}
\mvfrow{误解}{"我买了八只不同的基金，已经很分散了。"}
\mvfrow{事实}{老潘在武汉一家医院的设备科上班，2025 年陆续买了 8 只基金，名字分别叫科技创新、数字经济、高端制造、先进制造、智能汽车、新质生产力、成长精选、产业升级。他觉得自己布局得很周全。\\ 后来他把 8 只基金的前十大重仓拉到一张表上对齐，去重之后只剩 23 家公司；其中 6 家同时出现在 7 只基金里，合计占他总仓位的 \textbf{41\%}。2026 年 3 月这 6 家一起回调，他 8 只全绿，一只都没躲开。\\ \textbf{他不是买了 8 只基金，他是用 8 种方式买了同一个东西。}分散要看穿透之后的持仓，不看产品数量。}
\end{mythbox}

判断分散有没有效，有一个笨办法：回想最近一次市场大跌，你账户里有没有哪一栏是绿的旁边站着一栏红的。如果所有格子都是同一个颜色，你手上就只有一个资产，它只是被拆成了八个名字。

\section{每一类资产，在什么时候救你，在什么时候害你}

大类资产不该按"收益高低"排序，该按"它在什么天气下管用"排序。

\begin{tablewrap}
\begin{xltabular}{\linewidth}{>{\raggedright\arraybackslash}p{0.20\linewidth}XX}
\toprule
\thead{资产} & \thead{它在什么时候救你} & \thead{它在什么时候害你} \\
\midrule
\textbf{股票} & 经济扩张、企业盈利上行、通胀温和。长期分享增长的唯一入口 & 衰退、盈利下修、利率快速上行的时候。而且它跌起来不通知你 \\
\textbf{利率债（国债、政金债）} & 衰退和通缩恐慌里它涨，央行降息时它涨。提供现金流，也提供"别的东西都在跌"时的缓冲 & 通胀上行加加息，久期越长伤得越重。2022 年就是这么发生的 \\
\textbf{现金与类现金} & 所有人都缺钱的时候你有钱。它的价值不在利息，在于它是一张随时能出手的期权 & 通胀高的年份，它被安静地磨掉购买力。放太多是慢性失血 \\
\textbf{黄金与实物资产} & 货币信用被怀疑、地缘冲突、实际利率下行的时候 & 实际利率上行、风险偏好回暖的时候。它不生息，等待是有成本的 \\
\textbf{自住房} & 你得有地方住，这份居住价值天天在兑现。温和通胀下杠杆还能放大收益 & 你需要现金的时候。它卖不快、切不开，而且和你的工作往往在同一个城市里同生共死 \\
\bottomrule
\end{xltabular}
\end{tablewrap}

最后一栏那句得展开一句：房子的坏处不只是流动性差。你在杭州工作，房子也在杭州，如果杭州的主导产业出问题，你的收入和你的房价会一起往下走。第 19 章说的"不要重复押注同一个变量"，在这里同样成立，只是很少有人这么看自己的房子。

\section{60/40 的教训}

六成股票、四成债券，是过去三十年被写进无数教材和产品说明书的默认答案。它背后的逻辑是股债负相关：股票跌的时候债券涨，互相垫着。

2022 年，这套逻辑当场失效。

\begin{egbox}{2022 年发生了什么}
那一年标普 500 全收益 \textbf{−18.1\%}，美国综合债券指数 \textbf{−13.0\%}，经典 60/40 组合跌了 \textbf{16\%} 上下，是 1937 年以来最差的年份之一。

两条腿一起断，因为砍它们的是同一把刀。\textbf{通胀起来 → 央行加息 → 债券价格按久期下跌}（第 10 章那套算术），同时\textbf{折现率抬高 → 股票估值被压缩}（第 6 章那个除数）。同一个变量，两个方向的伤害。

持有 60/40 的人以为自己买了两样东西，2022 年才发现自己买的是同一个宏观赌注的两个版本。

\end{egbox}

这里要下一个判断，而且我认为它比大多数配置建议重要：\textbf{股债负相关不是自然法则，它是低通胀时代的产物。}

1998 到 2020 年那二十来年，冲击主要来自需求侧：需求塌了，企业盈利下滑，股票跌；央行降息救市，债券涨。负相关是这个因果链的副产品，不是资产本身的属性。

当冲击来自供给侧，链条反过来：能源涨价、关税加码、供应链断裂，物价上去了，央行不得不加息，股和债同向下跌。2026 年的世界里，供给侧的麻烦比过去二十年密集得多：中东局势推高油价并恶化贸易条件，关税与去全球化抬高了商品成本，AI 数据中心的资本开支在抢电力和产能。

在一个通胀中枢更高、供给冲击更频繁的世界里，只靠股债两条腿是不够的。这就是通胀对冲类资产在 2026 年被重新拿出来讨论的原因。

\section{央行在买黄金，但他们买的不是收益}

先把黄金的缺点说完，再说它为什么还在桌上。

黄金不产生现金流。一盎司金子放一百年，还是一盎司金子，不分红不生息。它的全部回报来自"下一个人愿意出多少钱"。长期看，它的实际回报低于股票，这一点没有争议。

\begin{mythbox}{常见误解}
\mvfrow{误解}{"黄金能保值，跌不到哪儿去。"}
\mvfrow{事实}{2011 年 9 月金价站上 1900 美元/盎司附近，到 2015 年 12 月跌到 1050 美元附近，\textbf{四年多跌掉四成半}，期间不产生一分钱利息。在 2011 年高点买入的人，要等到 2020 年才回到成本线。\\ 黄金对冲的是特定的风险：货币信用、地缘冲突、极端通胀。它不是"稳"的东西。把一件对冲工具当成保本工具，是 2013 年让很多人亏钱的那个误会。}
\end{mythbox}

说完缺点，说 2026 年最有说服力的那组证据。它不来自任何一位分析师的观点，来自各国央行的实际行为。

\begin{databox}{2026 数据：全球央行储备结构}
欧洲央行 2026 年 6 月《欧元的国际角色》报告：\textbf{截至 2025 年底，黄金已超越美国国债，成为全球央行的第一大储备资产。}

\begin{barfig}
  \barrow{黄金}{27}{27\%}
  \barrow{美国国债}{22}{22\%}
  \barrow{其他美元资产}{20}{20\%}
  \barrow{欧元}{15}{15\%}
\end{barfig}

黄金占比一年之间从 20\% 升到 27\%，美债从 25\% 降到 22\%，其他美元资产从 22\% 降到 20\%，欧元维持 15\%。美元计价资产合计仍占 \textbf{42\%}，仍是最大的单一币种，这一点没变。

\end{databox}

\begin{statrow}{3}
  \statitem{3.6}{万吨}{全球央行黄金储备，接近布雷顿森林体系鼎盛时期的约 3.8 万吨}
  \statitem{850}{吨}{2025 年全球央行净购金量}
  \statitem{5600}{美元}{2026 年 1 月金价一度接近的每盎司水平}
\end{statrow}

还有个有意思的细节：2025 年单一最大的黄金买家不是哪个国家，是稳定币发行商 Tether，一年买了超过 100 吨。一家发行数字美元的公司，把储备的一部分换成了地下室里的金属。

央行为什么买？\textbf{他们买的不是收益，是"不会被没收"。}一国的美债储备存在别人的清算系统里，别人可以按一个键把它冻住。这件事在 2022 年之后从假想题变成了现实题。金库里的实物金砖，放在自己国境内，谁也按不了那个键。

这个动机往下推一步，会推出一个对你有用的结论：\hlmark{央行的理由和你的理由不一样，所以他们的工具你不能照抄。}你在手机上买的黄金 ETF，登记在券商的账户里，依赖的恰恰是同一套清算系统。它能对冲通胀和货币贬值，但它对冲不了央行真正担心的那件事。你买它可以，但别把别人的动机当成自己的理由。

\hlwarm{这不构成投资建议。}本书不预测金价，2026 年 1 月的 5600 美元是一个已经发生的数字，不是一个方向。

\section{再平衡和定投：两条治人的规则}

配置比例定下来只是第一天的事。市场会替你改比例：涨得多的那一块自动变大，跌得多的自动缩小，一年以后你的 60/40 可能已经变成 71/29，而你什么都没做。

\term{再平衡}{rebalancing}\index{079@再平衡}就是把它掰回去。

\begin{flowlist}
  \flowitem{01}{按时间}{每年固定一天，全部调回目标比例。挑一个跟市场无关的日子，比如生日后的第一个工作日。别挑年底，年底所有人都在动，你的情绪会被新闻带着走。}
  \flowitem{02}{按阈值}{目标权益 60\%，涨到 65\% 或者跌到 55\% 就动手，中间不管。这个办法比按时间更贴合市场，代价是你得每个月看一眼。}
  \flowitem{03}{混合}{每季度看一次，只在偏离超过 5 个百分点时才操作。多数人用这个够了。}
  \flowitem{04}{用新钱补，别用卖出调}{每月的储蓄优先买进比例偏低的那一块。同样达到再平衡的效果，省掉申赎费和可能的税。这一条能省下的钱，往往比你纠结选哪只基金省下的多。}
\end{flowlist}

再平衡的价值不在于它聪明，在于它强迫你在下跌的时候买入、在上涨的时候卖出。这两件事人靠意志力做不到。每一轮下跌里基金赎回量都会放大，这是行业里公开的规律，说明大部分人在那个时刻做的恰好是相反的动作。

它是一台自动的高抛低吸机器，而且它不预测，只按比例干活。比喻可以继续推：什么时候这台机器会害你？当某个资产进入不可逆的单边下行时，机械再平衡会不断往一个正在归零的东西里加钱。所以\textbf{再平衡只能在大类资产层面做，不能在个股层面做}。一个国家的宽基指数不会归零，除非这个经济体消失；一家公司会。

说完再平衡说定投，这两件事经常被混为一谈，其实解决的是不同的问题。

\term{定投}{dollar-cost averaging}\index{012@定投}不提高预期收益。如果一个市场长期向上，一次性投入的期望收益高于分批投入，因为你的钱在场的时间更长。这个结论和很多人的直觉相反，但它就是算术。

定投降低的是另外两样东西：择时风险，你不会恰好把全部身家押在最高点；行为风险，市场跌的时候你还在买，因为扣款是自动的，不需要你每个月重新鼓一次勇气。

\begin{pullquote}{}
定投主要治的是人，不是市场。
\end{pullquote}

"微笑曲线"那套话术的问题在于它偷偷塞了个前提：市场一定会回来。

日经 225 在 1989 年 12 月见到 38915 点，之后花了 34 年，到 2024 年 2 月才第一次超越那个位置；2026 年 6 月首次突破 70000 点。从 1990 年开始定投日经的人，最后确实赢了，代价是熬三十四年，中间要眼睁睁看着账户长期浮亏并且每个月照扣不误。

所以定投真正的前提不是"曲线会微笑"，是两条：你买的那个东西背后的经济体还在增长，以及你真的能坚持三十年。第二条比第一条难得多。

\section{你的年龄，和中国投资者的三条绳子}

回到林悦。她 32 岁，人力资本 383 万，不含房的金融资产 25.6 万。她的"总资产"里 94\% 是人力资本。

那 25.6 万如果全部亏光，是她总资产的 6\%。同样一场跌 30\% 的熊市，落在她身上是总资产的 1.8\%，落在一个 62 岁、人力资本已经接近零的人身上，就是 30\%。

年轻人能承受更高的波动，不是因为胆子大，是因为后面还有几十年的现金流可以往坑里填。这就是\term{生命周期投资}{life-cycle investing}\index{047@生命周期投资}的全部道理。

但它有个前提，第 19 章那张职业表已经说过：如果你的人力资本本身就像一只高贝塔股票，这份额外的容忍度要打折。券商营业部里 28 岁的年轻人，未必比 45 岁的公务员能承受更多波动。

接下来是三条只在中国存在的绳子，绕不开。

\begin{egbox}{绳子一：房子已经把仓位占满了}
央行 2019 年城镇居民家庭资产负债调查显示，住房资产占城镇家庭总资产的 \textbf{59.1\%}。口径是 2019 年、城镇家庭，之后没有同口径的公开更新，但结构没有理由发生根本变化。

这意味着一件很尴尬的事：大多数中国家庭在打开基金 App 之前，已经默认满仓了一个加了杠杆、流动性极差、和自己收入高度相关的单一资产。

林悦的例子是现成的：房子 328 万，其余金融资产 25.6 万。她在那 25.6 万里怎么分配，对整体风险的影响不到一成。\textbf{对她来说，最有效的一次资产配置决策不在账户里，在于要不要买第二套。}

\end{egbox}

\textbf{绳子二：跨境配置是有额度的。}每人每年 5 万美元的便利化额度，QDII 基金的额度时紧时松、经常溢价，港股通的标的范围有限。教科书上那种"全球分散"的最优组合，在中国落地时会撞到一堵墙。承认这堵墙的存在，比假装它不在要有用。

\textbf{绳子三：你其实已经持有一大块债券了，只是它不叫债券。}林悦的公积金账户 9.6 万，加上养老保险个人账户里的余额，这些钱利率低、几乎没有波动、期限极长、强制缴纳。它的性质接近一笔久期很长的低息债券。

推论很实用：如果你已经被强制持有了这么一块类债券资产，你在金融账户里再配 60\% 的债，可能是在重复同一件事。至于该配多少，取决于你的公积金余额、你的年龄和你三年内要花的钱，这三个数只有你自己知道。

\begin{deepbox}{进阶：为什么没有通用比例}
本章从头到尾没有出现一个"建议配置比例"，这是故意的。\hlwarm{这不构成投资建议。}

原因不复杂：两个同样 35 岁、同样年入 40 万、在同一份风险测评问卷上得了同一档分数的人，一个两年后要付 90 万首付，一个未来十年没有任何刚性大额支出。他们能承受的波动差着十万八千里，而问卷把他们分进了同一格。

问卷问的是你的性格，配置该问的是你的日期表。任何人给你一个通用比例，都必须先问你两件事：这笔钱什么时候要用，以及你的刚性支出每月是多少。不问这两件事就给比例的，给的不是配置方案。

\end{deepbox}

\begin{takeaway}{本章一句话}
把涨跌不同步的东西放在一起，收益不减，波动变小。这是金融学唯一的免费午餐，而它的开关是相关性，不是数量。

\begin{itemize}
  \item 八只重仓同一批公司的基金，是一只基金的八个名字。分散看穿透持仓。
  \item 股债负相关是低通胀时代的产物，不是自然法则。2022 年两条腿一起断，因为砍它们的是同一把刀。
  \item 央行买黄金买的是"不会被没收"，你买的是账户里的一行数字。动机不同，别照抄。
  \item 再平衡强迫你反着人性做，定投治的是人不是市场。两者都不提高预期收益，都降低你搞砸的概率。
\end{itemize}

\end{takeaway}

\begin{bridgebox}
\textbf{下一章}配置好了，剩下的交给时间。可时间这件事，被误解得比什么都深。
\end{bridgebox}

